\documentclass[11pt,a4paper]{article}

% ================ Encodage et langue ================
\usepackage[utf8]{inputenc}
\usepackage[T1]{fontenc}
\usepackage[french]{babel}
\usepackage{lmodern}
\usepackage{textcomp}

% ================ Mise en page ================
\usepackage[margin=2.2cm]{geometry}
\usepackage{parskip}
\setlength{\parindent}{0pt}
\setlength{\parskip}{0.6em}

% ================ Couleurs et liens ================
\usepackage[table]{xcolor}
\definecolor{primary}{HTML}{0B5394}
\definecolor{accent}{HTML}{1F77B4}
\definecolor{light}{HTML}{F2F2F2}
\definecolor{codebg}{HTML}{F6F8FA}
\definecolor{rulec}{HTML}{D0D7DE}

\usepackage[colorlinks=true, linkcolor=primary, urlcolor=primary, citecolor=primary]{hyperref}

% ================ Tableaux ================
\usepackage{longtable, booktabs, array, tabularx}
\renewcommand{\arraystretch}{1.25}
\arrayrulecolor{rulec}

% ================ Listes ================
\usepackage{enumitem}
\setlist{itemsep=2pt, topsep=4pt}

% ================ Code et citations ================
\usepackage{fancyvrb}
\usepackage{framed}
\definecolor{shadecolor}{named}{light}

% Encadrés pour les blocs > (quotes)
\usepackage{mdframed}
\newmdenv[
  linecolor=primary,
  linewidth=2pt,
  topline=false, bottomline=false, rightline=false,
  innerleftmargin=10pt, innerrightmargin=8pt,
  innertopmargin=6pt, innerbottommargin=6pt,
  backgroundcolor=codebg,
  skipabove=8pt, skipbelow=8pt
]{quotebox}
\renewenvironment{quote}{\begin{quotebox}}{\end{quotebox}}

% ================ Titres ================
\usepackage{titlesec}
\titleformat{\section}{\Large\bfseries\color{primary}}{\thesection.}{0.6em}{}
\titleformat{\subsection}{\large\bfseries\color{accent}}{\thesubsection}{0.6em}{}
\titleformat{\subsubsection}{\normalsize\bfseries}{\thesubsubsection}{0.6em}{}

% ================ En-têtes et pieds de page ================
\usepackage{fancyhdr}
\pagestyle{fancy}
\fancyhf{}
\fancyhead[L]{\small\color{gray}ESPRIT School of Business --- M1 BA --- Projet Intégré}
\fancyhead[R]{\small\color{gray}\nouppercase{\leftmark}}
\fancyfoot[C]{\small\thepage}
\renewcommand{\headrulewidth}{0.3pt}
\renewcommand{\footrulewidth}{0pt}

% ================ Pandoc helpers ================
\providecommand{\tightlist}{\setlength{\itemsep}{0pt}\setlength{\parskip}{0pt}}
\providecommand{\passthrough}[1]{#1}

\usepackage{calc}
\providecommand{\real}[1]{#1}


% ================ Métadonnées ================
\title{\color{primary}\textbf{Description des Données}}
\author{Aymen Ben Brik}
\date{2026}

\begin{document}

\begin{titlepage}
\centering
\vspace*{1.5cm}

{\Large\bfseries ESPRIT School of Business}\\[0.3cm]
{\large Master 1 --- Business Analytics (M1 BA)}\\[2cm]

{\color{gray}\large Projet Intégré}\\[0.3cm]
{\color{gray}\large Analyse et Prédiction du Churn Client}\\[2.2cm]

\rule{0.85\linewidth}{1.2pt}\\[0.6cm]
{\Huge\bfseries\color{primary} Description des Données\par}
\rule{0.85\linewidth}{1.2pt}\\[2cm]

{\large Tuteur du projet}\\[0.2cm]
{\Large\bfseries Aymen Ben Brik}\\[1cm]
{\large 2026}

\vfill
{\color{gray}\small Document à destination des étudiants}
\end{titlepage}

\tableofcontents
\newpage

\section{Description des Données}\label{description-des-donnuxe9es}

Les données mises à votre disposition proviennent du système
d'information d'une institution bancaire. Elles ont été
\textbf{anonymisées} : aucun identifiant personnel direct n'est
exploitable, mais elles conservent leur structure et leur cohérence
métier réelle.

\subsection{1. Vue d'ensemble}\label{vue-densemble}

Le jeu de données est composé de :

\begin{itemize}
\tightlist
\item
  \textbf{1 fichier de données principal} au format CSV
  (\textasciitilde130 Mo, plusieurs centaines de milliers de lignes)
  contenant les informations clients, comptes et produits.
\item
  \textbf{8 tables de dimensions} au format Excel (\texttt{.xlsx}) qui
  donnent le libellé descriptif associé aux codes présents dans le
  fichier principal.
\end{itemize}

Chaque ligne du fichier principal correspond à un \textbf{couple
(client, compte produit)} : un même client peut donc apparaître
plusieurs fois s'il possède plusieurs comptes.

\subsection{2. Fichier principal --- Description des
colonnes}\label{fichier-principal-description-des-colonnes}

\subsubsection{2.1 Identifiants}\label{identifiants}

{\def\LTcaptype{none} % do not increment counter
\begin{longtable}[]{@{}lll@{}}
\toprule\noalign{}
Colonne & Type & Description \\
\midrule\noalign{}
\endhead
\bottomrule\noalign{}
\endlastfoot
\texttt{CUSTOMER\_NO} & Numérique & Identifiant unique du client \\
\texttt{ACCOUNT\_NO} & Numérique & Identifiant unique du compte \\
\texttt{BRANCH} & Code & Agence de rattachement \\
\end{longtable}
}

\subsubsection{2.2 Informations client}\label{informations-client}

{\def\LTcaptype{none} % do not increment counter
\begin{longtable}[]{@{}
  >{\raggedright\arraybackslash}p{(\linewidth - 4\tabcolsep) * \real{0.3333}}
  >{\raggedright\arraybackslash}p{(\linewidth - 4\tabcolsep) * \real{0.3333}}
  >{\raggedright\arraybackslash}p{(\linewidth - 4\tabcolsep) * \real{0.3333}}@{}}
\toprule\noalign{}
\begin{minipage}[b]{\linewidth}\raggedright
Colonne
\end{minipage} & \begin{minipage}[b]{\linewidth}\raggedright
Type
\end{minipage} & \begin{minipage}[b]{\linewidth}\raggedright
Description
\end{minipage} \\
\midrule\noalign{}
\endhead
\bottomrule\noalign{}
\endlastfoot
\texttt{NATIONALITY} & Code ISO & Nationalité du client \\
\texttt{RESIDENCE} & Code ISO & Pays de résidence \\
\texttt{MARITAL\_STATUS} & Code & Statut marital (C = célibataire, M =
marié(e), D = divorcé(e), V = veuf/veuve) \\
\texttt{DATE\_OF\_BIRTH} & Date (YYYYMMDD) & Date de naissance \\
\texttt{CUST\_OPENING\_DATE} & Date (YYYYMMDD) & Date d'entrée en
relation avec la banque \\
\texttt{NATURE\_CLIENT} & Code & Type de client (personne physique,
personne morale, etc.) \\
\texttt{PARTYCLASS} & Code & Classification commerciale (Retail,
Corporate, etc.) \\
\texttt{LOB} & Code & Ligne métier (Line of Business) \\
\texttt{SCORE\_KYC} & Code & Score de connaissance client (Know Your
Customer) \\
\texttt{COMPLETED\_FILE} & Booléen (YES/NO) & Dossier client complet \\
\texttt{LAST\_REVIEW\_DATE} & Date & Date de dernière revue du
dossier \\
\texttt{NEXT\_\_REVIEW\_DATE} & Date & Date de prochaine revue prévue \\
\texttt{INDUSTRY} & Code & Secteur d'activité du client (FK vers
\texttt{dim\_INDUSTRY}) \\
\texttt{SALARY} & Numérique & Salaire ou revenu déclaré \\
\end{longtable}
}

\subsubsection{2.3 Informations compte}\label{informations-compte}

{\def\LTcaptype{none} % do not increment counter
\begin{longtable}[]{@{}
  >{\raggedright\arraybackslash}p{(\linewidth - 4\tabcolsep) * \real{0.3333}}
  >{\raggedright\arraybackslash}p{(\linewidth - 4\tabcolsep) * \real{0.3333}}
  >{\raggedright\arraybackslash}p{(\linewidth - 4\tabcolsep) * \real{0.3333}}@{}}
\toprule\noalign{}
\begin{minipage}[b]{\linewidth}\raggedright
Colonne
\end{minipage} & \begin{minipage}[b]{\linewidth}\raggedright
Type
\end{minipage} & \begin{minipage}[b]{\linewidth}\raggedright
Description
\end{minipage} \\
\midrule\noalign{}
\endhead
\bottomrule\noalign{}
\endlastfoot
\texttt{ACCOUNT\_STATUS} & Texte & Statut du compte (\texttt{Active},
\texttt{Closed}, etc.) --- \textbf{variable cible pour le churn} \\
\texttt{ACCT\_OPENING\_DATE} & Date & Date d'ouverture du compte \\
\texttt{ACCT\_CLOSE\_DATE} & Date & Date de clôture du compte (si
applicable) \\
\texttt{ACCOUNT\_CATEGORY} & Code & Catégorie du compte (FK vers
\texttt{dim\_CATEGORY.ACCOUNT}) \\
\texttt{ACCOUNT\_TYPE\_DESC} & Texte & Description du type de compte \\
\texttt{CURRENCY} & Code ISO & Devise du compte (FK vers
\texttt{dim\_CURRENCY}) \\
\texttt{CLOSURE\_REASON} & Code & Motif de clôture (FK vers
\texttt{dim\_Closure\_reason}) \\
\texttt{ACCT\_BALANCE} & Numérique & Solde du compte \\
\end{longtable}
}

\subsubsection{2.4 Informations produit}\label{informations-produit}

{\def\LTcaptype{none} % do not increment counter
\begin{longtable}[]{@{}lll@{}}
\toprule\noalign{}
Colonne & Type & Description \\
\midrule\noalign{}
\endhead
\bottomrule\noalign{}
\endlastfoot
\texttt{PRODUCT\_GROUP} & Code & Groupe de produits \\
\texttt{PRODUCT\_LINE} & Code & Ligne de produits \\
\texttt{PRODUCT} & Code & Produit spécifique \\
\texttt{ACCOUNTNATURE} & Texte & Nature du produit \\
\texttt{STARTDATE} & Date & Date de début du produit \\
\texttt{MATURITYDATE} & Date & Date d'échéance \\
\texttt{AMOUNT} & Numérique & Montant du produit (crédit, dépôt,
etc.) \\
\texttt{FIXEDRATE} & Numérique & Taux fixe associé \\
\texttt{PRODUCT\_STATUS} & Texte & Statut du produit \\
\end{longtable}
}

\subsection{3. Tables de dimensions}\label{tables-de-dimensions}

Les fichiers \texttt{dim\_*.xlsx} fournissent les libellés descriptifs
associés aux codes du fichier principal. Elles sont à utiliser pour
enrichir les données lors de la phase ETL.

{\def\LTcaptype{none} % do not increment counter
\begin{longtable}[]{@{}
  >{\raggedright\arraybackslash}p{(\linewidth - 4\tabcolsep) * \real{0.3333}}
  >{\raggedright\arraybackslash}p{(\linewidth - 4\tabcolsep) * \real{0.3333}}
  >{\raggedright\arraybackslash}p{(\linewidth - 4\tabcolsep) * \real{0.3333}}@{}}
\toprule\noalign{}
\begin{minipage}[b]{\linewidth}\raggedright
Fichier
\end{minipage} & \begin{minipage}[b]{\linewidth}\raggedright
Clé
\end{minipage} & \begin{minipage}[b]{\linewidth}\raggedright
Contenu
\end{minipage} \\
\midrule\noalign{}
\endhead
\bottomrule\noalign{}
\endlastfoot
\texttt{dim\_CATEGORY.ACCOUNT.xlsx} & \texttt{CATEGORY\_id} &
Description des catégories de compte \\
\texttt{dim\_CURRENCY.xlsx} & Code devise & Description des devises \\
\texttt{dim\_Closure\_reason.xlsx} & \texttt{RECID} & Motifs de clôture
de compte \\
\texttt{dim\_DAO.xlsx} & Code & Référentiel DAO \\
\texttt{dim\_INDUSTRY.xlsx} & \texttt{INDUSTRY\_CODE} & Secteurs
d'activité économique \\
\texttt{dim\_SECTOR.xlsx} & Code & Secteurs détaillés \\
\texttt{dim\_TARGET.xlsx} & Code & Segmentation cible client \\
\texttt{dim\_TRANSACTION.xlsx} & Code & Référentiel de types de
transactions \\
\end{longtable}
}

\subsection{4. Définition de la variable cible
(churn)}\label{duxe9finition-de-la-variable-cible-churn}

Le \textbf{churn} se définit ici comme la \textbf{clôture d'un compte
client}. La variable cible peut être construite à partir de la colonne
\texttt{ACCOUNT\_STATUS} :

\begin{verbatim}
churn = 1  si  ACCOUNT_STATUS contient "Closed"
churn = 0  sinon
\end{verbatim}

Vous pouvez aussi proposer une définition plus fine (par exemple : churn
d'un client = clôture de \textbf{tous} ses comptes), à justifier dans le
rapport.

\subsection{5. Variables potentiellement intéressantes pour le
ML}\label{variables-potentiellement-intuxe9ressantes-pour-le-ml}

Quelques pistes (non exhaustives) à explorer :

\begin{itemize}
\tightlist
\item
  \textbf{Variables démographiques} : âge (à calculer depuis
  \texttt{DATE\_OF\_BIRTH}), statut marital, nationalité.
\item
  \textbf{Variables d'ancienneté} : ancienneté du client
  (\texttt{CUST\_OPENING\_DATE}), ancienneté du compte
  (\texttt{ACCT\_OPENING\_DATE}).
\item
  \textbf{Variables financières} : solde (\texttt{ACCT\_BALANCE}),
  salaire (\texttt{SALARY}), montant du produit (\texttt{AMOUNT}).
\item
  \textbf{Variables comportementales} : nombre de comptes par client,
  diversité des produits détenus, secteur d'activité, segment
  commercial.
\item
  \textbf{Variables de conformité} : score KYC, dossier complet ou non,
  ancienneté de la dernière revue.
\end{itemize}

\subsection{6. Qualité des données --- points
d'attention}\label{qualituxe9-des-donnuxe9es-points-dattention}

Le jeu de données est issu d'un système opérationnel réel. Vous
rencontrerez donc :

\begin{itemize}
\tightlist
\item
  Des \textbf{valeurs manquantes} dans plusieurs colonnes
  (\texttt{SALARY}, \texttt{INDUSTRY}, dates, etc.).
\item
  Des \textbf{valeurs \texttt{NULL} textuelles} à distinguer des
  \texttt{NaN}.
\item
  Des \textbf{dates au format \texttt{YYYYMMDD} numérique} à convertir.
\item
  Des \textbf{incohérences} possibles (dates de clôture sans statut
  \texttt{Closed}, soldes négatifs, etc.).
\item
  Un \textbf{déséquilibre de classes} entre comptes actifs et clos qu'il
  faudra prendre en compte dans le modèle ML.
\end{itemize}

Une phase de \textbf{profilage et d'analyse exploratoire} est donc
indispensable avant toute modélisation.

\subsection{7. Volumétrie indicative}\label{volumuxe9trie-indicative}

\begin{itemize}
\tightlist
\item
  Fichier principal : \textasciitilde130 Mo, plusieurs centaines de
  milliers de lignes.
\item
  Tables de dimensions : entre 6 ko et 80 ko chacune.
\end{itemize}

La volumétrie est suffisamment importante pour justifier une approche
structurée (Pandas avec chunks, ou base de données analytique), mais
reste manipulable sur une machine standard.

\subsection{8. Précautions}\label{pruxe9cautions}

\begin{itemize}
\tightlist
\item
  \textbf{Ne pas publier les données} sur GitHub ou tout autre dépôt
  public.
\item
  \textbf{Ne pas mentionner} l'origine des données dans aucun livrable.
\item
  Utiliser un \texttt{.gitignore} qui exclut au minimum :
  \texttt{*.csv}, \texttt{*.xlsx}, \texttt{*.pkl}, \texttt{data/}.
\end{itemize}

\end{document}
